\begin{abstract}

eBPF is widely adopted in modern production systems, yet its in-kernel JIT compilers remain under-optimized compared to mature runtimes like JVM and LLVM. The JIT intentionally trades performance for safety, and closing the gap inside the kernel would require tens of thousands of lines of safety-critical code. We find that even when optimizations are added, statically applying them regresses 68\% of real programs because optimization profitability is deployment-dependent. Our key insight is that eBPF bytecode, designed as a virtual ISA, is better understood as a \emph{deployment-time intermediate representation}: structured enough for optimization yet re-verifiable for safety, enabling an untrusted optimizer for kernel extensions. We present \tool, a deployment-time compilation framework that treats loaded bytecode as an IR and applies two classes of optimization passes: runtime specialization using deployment context, and hardware specialization through \emph{kinsn}, an extensible instruction mechanism analogous to LLVM intrinsics with formally verified equivalence between each eBPF expansion and its native lowering. We implement \tool with fewer than 600 lines of kernel modifications, transparently supporting the existing eBPF ecosystem. Evaluation demonstrates up to 40\% performance improvement, 50\% reduction in binary size, and low recompilation overhead. \tool provides a practical, upstream-compatible path toward extensible compilation frameworks for eBPF.

\end{abstract}

%%%%%%%%%%%%Outline%%%%%%%%%%%%%%%%%%


%%%%%%%%%%%%% PREVIOUS VERSION (kinsn-centric framing) — kept for comparison %%%%%%%%%%%%%
%
% \begin{abstract}
%
% Kernel extensions such as eBPF are widely used in performance and safety-critical paths. In the standard deployment path, programs are compiled to ordinary eBPF bytecode, verified in the kernel, and JIT-compiled to native code before execution. The same ordinary eBPF bytecode therefore serves both as the verifier's proof form and as the JIT's input to native code generation.
% %
% This bytecode form is good for verifier reasoning, but it erases optimization-relevant semantics before a program reaches the kernel. To preserve those semantics without complicating the in-kernel JIT, we need to separate the bytecode used for optimization from the bytecode used for verification. We propose \tool, a dynamic recompilation framework that realizes this separation.
% %
% \tool introduces kinsn, an extensible optimization-facing ISA whose instructions canonically expand to ordinary eBPF for kernel verification and lower directly to efficient native code. Optimization runs in userspace over kinsn, while the kernel reuses the existing verifier, re-JITs accepted rewrites, and installs them in place.
% %
% Our evaluation shows that \tool improves performance and reduces code size across diverse workloads with low recompilation overhead. We further demonstrate its effectiveness across architectures and its ability to preserve safety through fail-safe execution and security-oriented case studies.
%
% \end{abstract}
%
% STRENGTHS of this version:
% - "same ordinary eBPF bytecode serves both as the verifier's proof form and as the JIT's input" — crisp observation
% - "separate the bytecode used for optimization from the bytecode used for verification" — clean problem statement
% - kinsn "two realizations" framing is technically precise and memorable
% - more concise than the plan-based version
%
% WEAKNESSES of this version:
% - centers entirely on kinsn, but >50% of optimizations (map_inline, const_prop, DCE, wide_mem, etc.) don't use kinsn
% - does not mention the framework / extensibility story
% - does not mention Safety ≠ Correctness insight
% - does not mention runtime specialization or security hardening capabilities
%
%%%%%%%%%%%%% END PREVIOUS VERSION %%%%%%%%%%%%%

%%%
% message<background>
% kernel extensions like BPF are performance and safety critical

%%%
% message<motivation>
% kernel compiler is hard to optimize
% reason: tradeoff: implement optimization is heavy, kernel code is safety sensitive
% deeper reason: boundary is not defined

%%%
% message<idea>
% needs a better defined xx, need separation
% safety mechanism should be in kernel, optimization should be in userspace
% advantage: better performance, less complexity in each component, does not break safety

%%%
% message<contribution>
% framework: rewrite bpf program transparently in userspace, verify and recompile in kernel
% kinsn: an extensible instruction substrate (kinsn) for efficient code generation

%%%
% message<evaluation>
% performance: outperform baseline on various workloads
% code size: reduce binary size
% safety

%%%%%%%%%%%%Outline%%%%%%%%%%%%%%%%%%